\documentclass[10pt,a4paper]{amsart}
\usepackage{amsmath,amssymb,amsthm,graphicx}
% Retained minimal preamble for the Stacks Project chapter varieties.tex.
% Verbatim excerpt of the chapter preamble: Xy-pic setup (native preamble lines 14-18), the native theorem
% declarations with the mathematical macros (native preamble lines 154-217), and the project own no-op
% definitions of the reference/slogan/history display environments (native wrapper.tex lines 1-3) that the
% counted statement uses for its slogan. Copied mechanically; nothing added, removed or redefined.
% For commutative diagrams we use Xy-pic
\usepackage[all]{xy}

% We use 2cell for 2-commutative diagrams.
\xyoption{2cell}
% Theorem environments.
%
\theoremstyle{plain}
\newtheorem{theorem}[subsection]{Theorem}
\newtheorem{proposition}[subsection]{Proposition}
\newtheorem{lemma}[subsection]{Lemma}

\theoremstyle{definition}
\newtheorem{definition}[subsection]{Definition}
\newtheorem{example}[subsection]{Example}
\newtheorem{exercise}[subsection]{Exercise}
\newtheorem{situation}[subsection]{Situation}

\theoremstyle{remark}
\newtheorem{remark}[subsection]{Remark}
\newtheorem{remarks}[subsection]{Remarks}

\numberwithin{equation}{subsection}

% Macros
%
\def\lim{\mathop{\mathrm{lim}}\nolimits}
\def\colim{\mathop{\mathrm{colim}}\nolimits}
\def\Spec{\mathop{\mathrm{Spec}}}
\def\Hom{\mathop{\mathrm{Hom}}\nolimits}
\def\Ext{\mathop{\mathrm{Ext}}\nolimits}
\def\SheafHom{\mathop{\mathcal{H}\!\mathit{om}}\nolimits}
\def\SheafExt{\mathop{\mathcal{E}\!\mathit{xt}}\nolimits}
\def\Sch{\mathit{Sch}}
\def\Mor{\mathop{\mathrm{Mor}}\nolimits}
\def\Ob{\mathop{\mathrm{Ob}}\nolimits}
\def\Sh{\mathop{\mathit{Sh}}\nolimits}
\def\NL{\mathop{N\!L}\nolimits}
\def\CH{\mathop{\mathrm{CH}}\nolimits}
\def\proetale{{pro\text{-}\acute{e}tale}}
\def\etale{{\acute{e}tale}}
\def\QCoh{\mathit{QCoh}}
\def\Ker{\mathop{\mathrm{Ker}}}
\def\Im{\mathop{\mathrm{Im}}}
\def\Coker{\mathop{\mathrm{Coker}}}
\def\Coim{\mathop{\mathrm{Coim}}}

% Boxtimes
%
\DeclareMathSymbol{\boxtimes}{\mathbin}{AMSa}{"02}

%
% Macros for moduli stacks/spaces
%
\def\QCohstack{\mathcal{QC}\!\mathit{oh}}
\def\Cohstack{\mathcal{C}\!\mathit{oh}}
\def\Spacesstack{\mathcal{S}\!\mathit{paces}}
\def\Quotfunctor{\mathrm{Quot}}
\def\Hilbfunctor{\mathrm{Hilb}}
\def\Curvesstack{\mathcal{C}\!\mathit{urves}}
\def\Polarizedstack{\mathcal{P}\!\mathit{olarized}}
\def\Complexesstack{\mathcal{C}\!\mathit{omplexes}}
% \Pic is the operator that assigns to X its picard group, usage \Pic(X)
% \Picardstack_{X/B} denotes the Picard stack of X over B
% \Picardfunctor_{X/B} denotes the Picard functor of X over B
\def\Pic{\mathop{\mathrm{Pic}}\nolimits}
\def\Picardstack{\mathcal{P}\!\mathit{ic}}
\def\Picardfunctor{\mathrm{Pic}}
\def\Deformationcategory{\mathcal{D}\!\mathit{ef}}
\newenvironment{reference}{}{}
\newenvironment{slogan}{}{}
\newenvironment{history}{}{}
\usepackage{cleveref}
\crefname{originallemma}{original lemma}{original lemmas}
\Crefname{originallemma}{Original lemma}{Original lemmas}
\makeatletter
\newlabel{lemma-bijection-irreducible-components}{{038F}{source}{}{sourcepointer-1}{}}
\newlabel{lemma-bijection-irreducible-components@cref}{{[originallemma][0][]038F}{[1][0][]source}{}{}{}}
\makeatother
\crefname{originallemma}{original lemma}{original lemmas}
\Crefname{originallemma}{Original lemma}{Original lemmas}
\makeatletter
\newlabel{lemma-geometrically-irreducible-function-field}{{054Q}{source}{}{sourcepointer-2}{}}
\newlabel{lemma-geometrically-irreducible-function-field@cref}{{[originallemma][0][]054Q}{[1][0][]source}{}{}{}}
\makeatother
\title{Geometric irreducibility is equivalent to irreducibility of the self product}
\author{The Stacks Project Authors}
\date{stacks-project:varieties.tex:local-snapshot-20261011; GNU FDL 1.2 or later}
\begin{document}
\maketitle
\noindent\textit{Editorial note.} All delivered mathematical text is the author's own, copied from the pinned Stacks Project chapter file varieties.tex; no statement or proof step was rewritten, reordered or removed.\par\medskip
\noindent\textit{Editorial note.} The proof is complete as written and ends with the function field criterion; no step is referred to the reader.\par\medskip
\noindent\textit{Editorial note.} The counted spans carry 2 navigation stubs and no bibliography citation.\par\medskip
\noindent\textit{Editorial note.} The preamble is the chapter's retained minimal preamble (sources/stacks-varieties/retained-preamble.tex lines 1-77), a verbatim excerpt of the native preamble.\par\medskip
\noindent\textit{Editorial note.} The label lemma-geometrically-irreducible-self-product is defined only in varieties.tex, so no cross-chapter restatement had to be weighed for this card.\par\medskip
\noindent\textit{Allowed standard prerequisites (editorial).}\begin{itemize}
\item Behaviour of irreducible components under the self product and base change. Same chapter, Lemma lemma-bijection-irreducible-components, invoked at proof lines 1493-1495. Supplies the easy direction from geometric irreducibility to irreducibility of the self product.
\item Criterion for geometric irreducibility through the function field. Same chapter, Lemma lemma-geometrically-irreducible-function-field, invoked at proof lines 1519-1520. Supplies that having no separably algebraic elements of the function field outside k implies geometric irreducibility.
\item A finite separable extension E of k has E tensor\_k E a product of fields of dimension dim\_k(E). Standard field theory, used at proof lines 1513-1517. Supplies the contradictory decomposition of E tensor\_k E into at least two fields.
\end{itemize}\medskip
\ifdefined\hypertarget\hypertarget{sourcepointer-1}{}\fi
\noindent\textit{Source reference.} Original lemma 038F: sources/stacks-varieties/pinned-original-source.tex, source lines 1380--1380. Pointer outside the counted span, invoked once for context; the load-bearing argument of the counted statement is the author's proof reproduced in full in the proof span.\par\medskip
\ifdefined\hypertarget\hypertarget{sourcepointer-2}{}\fi
\noindent\textit{Source reference.} Original lemma 054Q: sources/stacks-varieties/pinned-original-source.tex, source lines 1446--1446. Pointer outside the counted span, invoked once for context; the load-bearing argument of the counted statement is the author's proof reproduced in full in the proof span.\par\medskip
\begin{lemma}
\label{lemma-geometrically-irreducible-self-product}
A scheme $X$ over a field $k$ is geometrically irreducible
over $k$ if and only if $X\times_k X$ is irreducible.
\end{lemma}

\begin{proof}
If $X$ is geometrically irreducible over $k$,
then $X\times_k X$ is irreducible
by Lemma \ref{lemma-bijection-irreducible-components}.

\medskip\noindent
Conversely, suppose that $X \times_k X$ is irreducible. Replacing $X$
by the underlying reduced scheme, we can assume that $X$ is reduced
without changing the problem.
Since the first projection $X \times_k X \to X$
is surjective, $X$ is also irreducible, hence integral.
Suppose that there is an element
$\alpha$ in the function field $k(X)$ that is separably
algebraic over $k$ but not in $k$. Let $E = k(\alpha)$.

\medskip\noindent
Then there is a nonempty affine open subscheme $U$ of $X$ such that
$k \subset E \subset O(U)$, and so $E \otimes_k E \subset
O(U)\otimes_k O(U) = O(U\times_k U)$. Thus we have a dominant morphism
$U \times_k U \to \Spec(E\otimes_k E)$. Since $X\times_k X$ is irreducible,
so is $U \times_k U$, and hence $\Spec(E \otimes_k E)$ is connected.
If $E$ is strictly bigger than $k$, then $E \otimes_k E$
is a product of fields including $E$ as one factor. It has dimension
$(\dim_k(E))^2>\dim_k(E)$ as a $k$-vector space,
and so $E\otimes_k E$ is a product
of at least two fields, contradicting that $\Spec(E \otimes_k E)$
is connected. So in fact $E$ is equal to $k$.
By Lemma \ref{lemma-geometrically-irreducible-function-field},
$X$ is geometrically irreducible over $k$.
\end{proof}

\end{document}
